% TOP_PROBLEM: {"id": 2295, "title": "Erdős Problem #769", "category_id": 1, "category": {"id": 1, "name": "number_theory", "display_name": "Number Theory", "description": "Properties of integers, prime numbers, Diophantine equations.", "slug": "number-theory", "order_index": 1, "created_at": "2026-07-31T15:26:25.670Z"}, "status": "open"}
% FIRST_POSTED: null
% ENTRY_KIND: statement_only
% Imported from: unsolvedmath_additional_500.tex; UnsolvedMath ID 2295
% !TeX program = lualatex
% Compile twice: lualatex 2295.tex
% Self-contained: no external bibliography, images, or \input files.
% Numeric section labels and visible numbers are original UnsolvedMath IDs.
\documentclass[11pt,a4paper]{article}
\usepackage[margin=24mm,headheight=14pt]{geometry}
\usepackage{fontspec}
\setmainfont{Latin Modern Roman}
\setsansfont{Latin Modern Sans}
\setmonofont{Latin Modern Mono}
\usepackage{amsmath,amssymb,mathtools}
\usepackage{unicode-math}
\setmathfont{Latin Modern Math}
\usepackage{microtype}
\usepackage{enumitem,xurl,needspace}
\usepackage[dvipsnames]{xcolor}
\usepackage{hyperref}
\hypersetup{unicode=true,colorlinks=true,linkcolor=MidnightBlue,urlcolor=MidnightBlue,
 pdftitle={UnsolvedMath Problem 2295},pdfauthor={Source-annotated compilation},bookmarksnumbered=true}
\usepackage{bookmark,tocloft,titlesec,fancyhdr}
\setlength{\cftsecnumwidth}{4.2em}
\cftsetpnumwidth{2.5em}
\cftsetrmarg{4em}
\titleformat{\section}{\Large\bfseries\raggedright}{\thesection}{0.7em}{}
\pagestyle{fancy}\fancyhf{}
\fancyhead[L]{\small\sffamily Additional UnsolvedMath statements}
\fancyhead[R]{\small Original catalogue IDs}
\fancyfoot[C]{\thepage}
\setlength{\parindent}{0pt}
\setlength{\parskip}{5pt plus 1pt}
\setlength{\emergencystretch}{3em}
\setcounter{tocdepth}{1}\setcounter{secnumdepth}{1}
\setlist[itemize]{leftmargin=3em,itemsep=4pt,topsep=4pt,parsep=0pt}
\allowdisplaybreaks
\newcommand{\N}{\mathbb{N}}
\newcommand{\Z}{\mathbb{Z}}
\newcommand{\Q}{\mathbb{Q}}
\newcommand{\R}{\mathbb{R}}
\newcommand{\C}{\mathbb{C}}
\newcommand{\F}{\mathbb{F}}
\newcommand{\A}{\mathbb{A}}
\newcommand{\PP}{\mathbb{P}}
\newcommand{\E}{\mathbb{E}}
\newcommand{\eps}{\varepsilon}
\newcommand{\dd}{\,\mathrm d}
\newcommand{\sref}[2]{\hyperlink{source:#1:#2}{[S#2]}}
\newif\ifshowcatalogue
\showcataloguetrue
\newcommand{\cataloguescope}[1]{\ifshowcatalogue
 \paragraph{Statement recorded in the supplied source.}
 #1\fi}
\begin{document}
% UnsolvedMath ID 2295; source catalogue code EP-769

\setcounter{section}{2294}

\section{Thresholds for Tiling Cubes with Prescribed Numbers of Cubes}\label{2295}

{\small\sffamily UnsolvedMath ID 2295 · Number Theory}

\subsection{Definitions and mathematical statement}

\paragraph{Shared notation.}Unless a section says otherwise, $\N=\{1,2,\ldots\}$,
$\Z$ is the ring of integers, $\Q,\R,\C$ are the rational, real, and complex
fields, $[N]=\{1,\ldots,\lfloor N\rfloor\}$, and $\log$ is the natural
logarithm. For real functions of a parameter tending to infinity,
$f\ll g$ and $f=O(g)$ mean $|f|\leq Cg$ eventually for some fixed $C$;
$f\gg g$ reverses this comparison; $f=o(g)$ means $f/g\to0$;
$f\sim g$ means $f/g\to1$; and $f\asymp g$ means both order bounds.
A subscript on $O$ or $\ll$ specifies allowed constant dependence.
The expression $n^{o(1)}$ in an upper bound means growth smaller than
$n^\epsilon$ for each fixed $\epsilon>0$. Local definitions govern any
exception, finite-field convention, or use of the same symbol for another
object. An asymptotic claim is not automatically an effective algorithm.



Homothetic cubes are translates and positive dilates of the original cube, with the same orientation. A decomposition has pairwise disjoint interiors and union equal to the unit cube. The threshold $c(n)$ must allow every larger integer tile count, not just a sequence of counts. \sref{2295}{1} \sref{2295}{2}

\paragraph{Statement recorded in the supplied source.}
Let $c(n)$ be minimal such that if $k\geq c(n)$ then the $n$-dimensional unit cube can be decomposed into $k$ homothetic $n$-dimensional cubes. Give good bounds for $c(n)$ - in particular, is it true that $c(n) \gg n^n$? \sref{2295}{1}

\paragraph{Residual task reported by the archive.}

The embedded review (2026-08-17) records: Prove or refute a uniform lower bound of order n\textasciicircum{}n and sharpen c(n), especially when n+1 is prime. \sref{2295}{1}

\subsection{Short English statement}

How large must a tile count be before every larger count can occur in a tiling of a fixed-dimensional cube by smaller homothetic cubes?

\subsection{Sources}

{\small

\begin{itemize}

\item[\hypertarget{source:2295:1}{S1}] ULAM.AI, \emph{UnsolvedMath}, supplied \texttt{problems.json}, record 2295 (EP-769), statement and background. The embedded literature-review date is 2026-08-17; that is the archive’s review date, not a fresh certification. Dataset landing page: \url{https://huggingface.co/datasets/ulamai/UnsolvedMath}.

\item[\hypertarget{source:2295:2}{S2}] T. F. Bloom, \emph{Erdős Problems}, Problem 769, \url{https://www.erdosproblems.com/769}. Reference imported from the supplied record; the live problem page was not independently checked for this compilation.

\item[\hypertarget{source:2295:3}{S3}] Peter Connor and Phillip Marmorino, Decomposing cubes into smaller cubes, Journal of Geometry 109 (2018), article 19. \url{https://doi.org/10.1007/s00022-018-0424-4}. Bibliographic information imported from the supplied record.

\item[\hypertarget{source:2295:4}{S4}] Connor, Peter and Marmorino, Phillip, Decomposing cubes into smaller cubes. J. Geom. (2018), Paper No. 19, 11. Bibliographic information imported from the supplied record.

\item[\hypertarget{source:2295:5}{S5}] Erd\H{o}s, P., Remarks on some problems in number theory. Math. Balkanica (1974), 197-202. Bibliographic information imported from the supplied record.

\item[\hypertarget{source:2295:6}{S6}] Hudelson, Matthew, Dissecting {$d$}-cubes into smaller {$d$}-cubes. J. Combin. Theory Ser. A (1998), 190--200. Bibliographic information imported from the supplied record.

\end{itemize}

}

\label{end:2295}
\end{document}
